% !TEX root = ../main.tex
\section{结论}

本文研究了一个面向复杂系统的基础逆问题：当可获得的信息主要是各组分的状态时间序列时，如何学习网络上的演化规律，并在连接关系不可见时进一步重构潜在拓扑。
围绕这一目标，
本文形成了由已知网络到未知网络的递进方法。
前者以自身--邻居双通道分别表示组分内在变化与局部交互，
后者在此基础上加入离散拓扑生成器，
通过硬图规律学习与软图结构更新交替优化。

网络 SIR 为这套方法提供了可解释的验证环境。
已知网络实验表明，分离表示能够缓解邻居累积量级与自身状态之间的干扰，在感染曲线、参数反演和多步推演上取得更稳定的综合结果。
未知网络实验进一步表明，
在 ER、WS 和 BA 三个固定实例上，
模型可以仅依据状态转移数据恢复 F1 为 0.937--1.000 的硬图，
并保持接近真实传播参数的演化响应和较低的自回归误差。

本文还建立了连接一步预测、规律诊断、多步推演和边级恢复的评价方式。
各层面结果相互印证，使模型性能既能以预测误差衡量，也能由局部机制和潜在连接的恢复加以解释，
为复杂系统的数据驱动建模提供了更完整的评价依据。

本文并不把 SIR 视为方法的边界。
自身演化与交互作用的分解、离散结构与动力学的联合约束，以及多层次评价原则，都具有迁移到其他网络动力学的潜力。
未来将通过更多动力学类型、独立网络实例、噪声和缺失观测以及有向、加权和时变网络，
系统检验这一潜力，
并推动方法由受控验证走向真实复杂系统的建模与预测。
